\section{Actividades y cargas de los núcleos funcionales}
Los 90 paquetes se derivan en 360 actividades por aporte de perfil, identificadas en \path{componentes/actividades_nucleos.csv}. Cada actividad conserva cuenta, resultado, ejecutor, responsable último, horas, predecesoras, lugar y ventana quincenal. Las comprobaciones de Calidad dependen de los aportes del paquete, y las vistas de la lógica/datos y los canales comunes. Su suma reproduce las 4.200 HH sin agregar gestión, pruebas unitarias o documentación ya incluidas en el paquete.

La especificación E1 toma la línea base de H1 y se contrasta con modelo, seguridad y arquitectura de H2 antes de construcción. BASE-E1 significa interfaces y decisiones de la primera etapa estabilizadas para diseñar E2 durante su marcha blanca; no exige esperar producción del mes 16 ni presume E1 recibida. La construcción E2 depende de su diseño detallado H8. El código nuevo conserva el inicio de construcción desde el mes 8 de SD1. Las ventanas iniciales de 8--9 para E1 y 15--16 para E2 permiten detectar recursos y enlaces necesarios, pero requieren nivelación antes de incorporarse a la línea base \parencite[art.~17; Formulario~E-25; pp.~12--74]{bases-admin}.

Las unidades y contratos se prueban dentro de construcción; el sistema completo se recibirá con UAT, cobertura de lógica, carga a 1,5 veces peak, estrés, seguridad ofensiva, accesibilidad, resiliencia y escenarios desconectados de 1.10. Los umbrales de tareas de SD3 se conservan como criterios de diseño y se demostrarán mediante esos protocolos. Ni terminar una ficha ni ejecutar una prueba unitaria acredita un tiempo de atención, ocho horas de inspección, doce horas de terminal o veinticuatro horas de célula local.

\begin{longtable}{R{12mm}R{18mm}R{18mm}R{18mm}R{18mm}R{18mm}R{18mm}R{18mm}}
\caption{Cargas iniciales de los núcleos funcionales}\label{tab:sd7-nucleos-curva}\\\hline
 \EncabezadoTabla{Mes} & \EncabezadoTabla{A} & \EncabezadoTabla{D} & \EncabezadoTabla{S} & \EncabezadoTabla{Q} & \EncabezadoTabla{V} & \EncabezadoTabla{I} & \EncabezadoTabla{HH}\\\hline\endfirsthead
\hline \EncabezadoTabla{Mes} & \EncabezadoTabla{A} & \EncabezadoTabla{D} & \EncabezadoTabla{S} & \EncabezadoTabla{Q} & \EncabezadoTabla{V} & \EncabezadoTabla{I} & \EncabezadoTabla{HH}\\\hline\endhead
3 & 288 & 144 & 72 & 0 & 0 & 72 & 576 \\\hline
8 & 0 & 144 & 0 & 72 & 864 & 72 & 1152\\\hline
9 & 72 & 0 & 72 & 72 & 576 & 0 & 792 \\\hline
13 & 192 & 96 & 48 & 0 & 0 & 48 & 384\\\hline
15 & 0 & 96 & 0 & 48 & 576 & 48 & 768 \\\hline
16 & 48 & 0 & 48 & 48 & 384 & 0 & 528\\\hline
\end{longtable}\FuenteTabla{Ventanas derivadas del reparto E1/E2; distribución inicial sin nivelación aprobada. No implica disponibilidad por la coincidencia de paquetes en un mes.}


\subsection{Balance conjunto y necesidad de revisar dotación}
\begin{longtable}{R{18mm}L{36mm}R{30mm}R{30mm}}
\caption{Balance parcial con los núcleos funcionales}\label{tab:sd7-nucleos-balance}\\\hline
Mes & Perfil/persona & HH conocidas & Exceso sobre 144\\\hline\endfirsthead
\hline Mes & Perfil/persona & HH conocidas & Exceso sobre 144\\\hline\endhead
3 & A & 344 & 200 \\\hline
3 & D & 248 & 104 \\\hline
6 & O+L & 160 & 16 \\\hline
8 & D & 304 & 160 \\\hline
8 & Q & 184 & 40 \\\hline
8 & V & 1208 & 1064 \\\hline
9 & V & 840 & 696 \\\hline
13 & A & 224 & 80 \\\hline
13 & D & 152 & 8 \\\hline
15 & V & 640 & 496 \\\hline
16 & Q & 180 & 36 \\\hline
16 & V & 656 & 512 \\\hline
17 & Q & 156 & 12 \\\hline
17 & V & 384 & 240 \\\hline
18 & Q & 168 & 24 \\\hline
20 & Q & 152 & 8 \\\hline
21 & O+L & 164 & 20 \\\hline
\end{longtable}\FuenteTabla{Suma de cuentas cuantificadas, incluido seguimiento fijo de innovación; excluye gobierno, pruebas/implantación general, servicio base y trabajo variable.}


Este balance integra las cuentas cuantificadas y sustituye para su alcance los contrastes parciales anteriores. Las ventanas concentradas muestran sobrecargas; no se utilizan como una promesa de fechas viable. Entre los meses 8 y 9 la carga conocida de Desarrollo suma $1.208+840=2.048$ HH. Una persona aporta $2\times144=288$ HH bajo la hipótesis de capacidad del T-15; mantener íntegramente esa ventana exigiría una cota inferior de $\lceil2.048/288\rceil=8$ personas de Desarrollo, antes de considerar secuencia, ausencias o correcciones. No se presume que esas personas estén contratadas ni que los ocho titulares de SD1 puedan ejercer como desarrolladores.

Incluso extendiendo toda esa carga hasta el mes 10, la cota sería $\lceil2.048/(3\times144)\rceil=5$ personas, sin demostrar que el retraso preserve ensayos, revisiones y H4/H5. No constituye una alternativa aprobada. La dependencia de versión integrada antes de los ensayos exige calcular qué trabajo puede anticiparse, qué parte puede avanzar en paralelo y qué refuerzo especializado deberá reservarse con nombre de función, disponibilidad y costo. El problema no se resuelve distribuyendo las mismas horas entre más meses si las dependencias o las reservas contractuales lo impiden.

Los 4.200 HH de núcleos elevan la estimación parcial de implementación a 11.060 HH, más revisión documental variable, enrolamiento y las cuentas de gobierno, pruebas integrales, capacitación/implantación general y salida aún por dimensionar. El subtotal de operación mantiene su separación. La Dirección completará esas cuentas y reconciliará la dotación, las dedicaciones de SD1 y el cronograma de SD13 antes de aprobar la línea base. La suma de fichas no demuestra todavía la regla de cobertura del 100\,\%.
